\section{Methodology}
\label{sec:method}

In this section, we present the formal mathematical formulation of \textbf{PAC-Bayes Merge}. We first describe the parameterization of ensembling coefficients via continuous polynomial trajectories in Section~\ref{subsec:poly_trajectory}. We then derive the information-theoretic PAC-Bayesian generalization bound under a spherical Gaussian prior in Section~\ref{subsec:pac_bayes_derivation}, proving that minimizing this bound analytically justifies a quadratic $L_2$ Consensus-Pulling penalty centered at the stable uniform baseline. In Section~\ref{subsec:initialization_consensus}, we determine the exact mathematical targets for initialization. In Section~\ref{subsec:swa_equivalence}, we establish the learning-theoretic equivalence between uniform weight merging and Stochastic Weight Averaging (SWA). Finally, in Section~\ref{subsec:monte_carlo}, we present the Monte Carlo expected risk optimization and posterior ensemble evaluation that bridges the theoretical-to-empirical gap.

\subsection{Weight-Space Merging and Polynomial Trajectory Parameterization}
\label{subsec:poly_trajectory}
Let $W_0^{(l)} \in \mathbb{R}^{D_l}$ represent the parameter weight tensor of a pre-trained base model at network layer $l \in \{0, \dots, L-1\}$, where $L$ is the total number of layers. Suppose we are given $K$ distinct task-specific expert models, each independently fine-tuned from the same shared pre-trained base $W_0$ on task $k \in \{0, \dots, K-1\}$. We define the corresponding parameter-space \emph{task vectors} at layer $l$ as:
\begin{equation}
V_k^{(l)} = W_k^{(l)} - W_0^{(l)}
\end{equation}

To merge these expert models into a single, cohesive multi-task parameter set, we define layer-wise ensembling coefficients $\alpha_k(l) \in [0, 1]$, representing the relative contribution of expert $k$ at layer $l$. Rather than optimizing independent, unconstrained coefficients at each layer---which creates a high-dimensional search space of size $K \times L$ and exposes the model to transductive overfitting---we restrict the coefficients to follow a global, continuous trajectory across network depth. 

Specifically, we normalize the network layer index to a continuous depth coordinate $z = \frac{l}{L-1} \in [0, 1]$. We parameterize the inner ensembling logits using a low-degree polynomial $p_k(z; \theta_k)$ of degree $d$ (typically $d \le 2$):
\begin{equation}
p_k(z; \theta_k) = \sum_{j=0}^d \theta_{k,j} z^j
\end{equation}
where $\theta_k = (\theta_{k,0}, \theta_{k,1}, \dots, \theta_{k,d})^T \in \mathbb{R}^{d+1}$ are the continuous polynomial coefficients for task expert $k$, and $\Theta = [\theta_0, \dots, \theta_{K-1}]^T \in \mathbb{R}^{K \times (d+1)}$ represents the complete learnable coordinate matrix. To bound each ensembling coefficient strictly to the interval $[0, 1]$, we apply the logistic sigmoid function $\sigma(u) = (1 + e^{-u})^{-1}$:
\begin{equation}
\alpha_k(l; \theta_k) = \sigma\left( p_k\left(\frac{l}{L-1}; \theta_k\right) \right) = \sigma\left( \sum_{j=0}^d \theta_{k,j} \left(\frac{l}{L-1}\right)^j \right)
\end{equation}
While a task-wise \texttt{softmax} function would constrain ensembling coefficients to the probability simplex $\Delta^{K-1} = \{x \in \mathbb{R}^K \mid x_i \ge 0, \sum_i x_i = 1\}$, doing so prevents the model from independently adjusting the scale of task vectors in intermediate layers where multi-task representation sharing requires non-convex combinations (with $\sum_k \alpha_k(l) \neq 1$). Our independent coordinate-wise parameterization provides a richer, non-convex parameter-space combination that preserves continuous representative capacity while being regularized towards the uniform consensus basin.

The merged weight tensor $W_{\text{merged}}^{(l)}(\Theta)$ at layer $l$ is then assembled by adding the weighted task vectors to the pre-trained base model:
\begin{equation}
W_{\text{merged}}^{(l)}(\Theta) = W_0^{(l)} + \sum_{k=0}^{K-1} \alpha_k(l; \theta_k) V_k^{(l)}
\end{equation}

By projecting the $K \times L$ ensembling parameters onto a low-degree polynomial space of size $K \times (d+1)$, this parameterization serves as a continuous, depth-wise parametric low-pass filter. It suppresses high-frequency coordinate oscillations across adjacent layers, preventing the chaotic coefficient transitions that characterize unconstrained optimization on few-shot calibration data.

From a learning-theoretic perspective, restricting the trajectory to a low-degree polynomial ($d \le 2$) serves as a structural capacity-bounding mechanism. By mapping the parameter space from a high-dimensional space of $K \times L = 56$ parameters to a compact subspace of dimension $K \times (d+1) = 12$, we severely constrain the complexity of the hypothesis class. Mathematically, this reduction in parameter-space dimensionality directly minimizes the Kullback-Leibler divergence $D_{\text{KL}}(Q \parallel P)$ in the PAC-Bayesian bound (as shown in Section~\ref{subsec:pac_bayes_derivation}), leading to a tighter and more stable generalization guarantee. This provides a formal learning-theoretic explanation for why low-degree polynomials are not just an empirical heuristic, but an optimal capacity control strategy for few-shot adaptation.

\subsection{PAC-Bayesian Generalization Derivation}
\label{subsec:pac_bayes_derivation}
While polynomial parameterization reduces the dimensionality of the search space, optimizing $\Theta$ on extremely scarce calibration data (e.g., $M = 10$ samples per task) is still vulnerable to transductive overfitting. To provide a rigorous learning-theoretic defense, we establish a formal PAC-Bayesian framework over our trajectory coordinate space.

Let $\mathcal{D}_{\text{cal}} = \cup_{k=0}^{K-1} \{(x_i^{(k)}, y_i^{(k)})\}_{i=1}^{M}$ be a multi-task calibration dataset containing $M$ samples per task across $K$ tasks, yielding a total calibration size of $N_{\text{total}} = K \times M$ (for our setup, $M = 10$ and $K = 4$, so $N_{\text{total}} = 40$). We define a randomized posterior classifier $G_{\tilde{\Theta}}$ where the trajectory parameters $\tilde{\Theta}$ are random variables drawn from a posterior probability distribution $Q$. We specify $Q$ to be an isotropic Gaussian distribution centered at our learnable trajectory coefficients $\Theta$, with a fixed posterior variance $\sigma^2 I$:
\begin{equation}
Q(\tilde{\Theta}) \sim \mathcal{N}\left(\Theta, \sigma^2 I_{K(d+1)}\right)
\end{equation}

Before observing the calibration data, we define a spherical Gaussian prior distribution $P$ centered at the stable, scale-preserving uniform ensembling consensus baseline $\Theta_{\text{uniform}}$, with a prior variance $\sigma_0^2 I$:
\begin{equation}
P(\tilde{\Theta}) \sim \mathcal{N}\left(\Theta_{\text{uniform}}, \sigma_0^2 I_{K(d+1)}\right)
\end{equation}

\begin{remark}
We note that specifying a fixed isotropic variance ($\sigma^2 I$ and $\sigma_0^2 I$) is mathematically convenient as it enables an analytical closed-form solution for the Kullback-Leibler divergence. However, in physical neural networks, parameter sensitivities are highly heterogeneous across layers (e.g., intermediate representations versus classifier heads). While treating the prior-posterior distributions as isotropic simplifies the PAC-Bayesian bound, future work could investigate layer-wise adaptive variances based on the empirical Fisher Information Matrix, as discussed in Section~\ref{sec:conclusion} and derived in Appendix~\ref{app:non_isotropic_derivation}.
\end{remark}

According to McAllester's classical PAC-Bayesian theorem~\cite{mcallester1999pac}, for any confidence level $\delta \in (0, 1)$, with probability at least $1-\delta$ over the random draw of the calibration set, the true out-of-distribution multi-task classification risk $R(G_Q) = \mathbb{E}_{\tilde{\Theta} \sim Q}[R(G_{\tilde{\Theta}})]$ is bounded by:
\begin{equation}
\label{eq:mcallester}
\mathbb{E}_{\tilde{\Theta} \sim Q}[R(G_{\tilde{\Theta}})] \le \mathbb{E}_{\tilde{\Theta} \sim Q}[\widehat{R}_{\text{cal}}(G_{\tilde{\Theta}})] + \sqrt{\frac{D_{\text{KL}}(Q \parallel P) + \ln\left(\frac{2\sqrt{N_{\text{total}}}}{\delta}\right)}{2 N_{\text{total}}}}
\end{equation}
where $\widehat{R}_{\text{cal}}$ represents the empirical classification risk (cross-entropy loss) on the calibration dataset $\mathcal{D}_{\text{cal}}$, and $D_{\text{KL}}(Q \parallel P)$ is the Kullback-Leibler divergence between the posterior and prior distributions.

Since $Q$ and $P$ are isotropic Gaussian distributions in $\mathbb{R}^{K(d+1)}$, their KL divergence can be computed analytically in closed form:
\begin{equation}
D_{\text{KL}}(Q \parallel P) = \int_{\mathbb{R}^{K(d+1)}} Q(\tilde{\Theta}) \ln \frac{Q(\tilde{\Theta})}{P(\tilde{\Theta})} d\tilde{\Theta}
\end{equation}
Substituting the probability density functions of the two multivariate Gaussians yields:
\begin{equation}
\label{eq:kl_derivation}
D_{\text{KL}}(Q \parallel P) = \frac{\|\Theta - \Theta_{\text{uniform}}\|_2^2}{2 \sigma_0^2} + \frac{K(d+1)}{2} \left( \frac{\sigma^2}{\sigma_0^2} - 1 - \ln \frac{\sigma^2}{\sigma_0^2} \right)
\end{equation}

Since the posterior variance $\sigma^2$ and prior variance $\sigma_0^2$ are fixed constants, the second term in Equation~\eqref{eq:kl_derivation} is a constant independent of the learnable trajectory mean $\Theta$. Because the generalization bound in Equation~\eqref{eq:mcallester} contains the Kullback-Leibler divergence inside a non-linear square root, direct optimization of this exact bound is highly non-convex. Specifically, since the KL divergence is quadratic in $\Theta$, the complexity penalty in McAllester's bound is actually linear in $\|\Theta - \Theta_{\text{uniform}}\|_2$ due to the square root.

Following standard practice in PAC-Bayesian literature (e.g., \cite{catoni2007pac, alquier2016properties}), direct optimization of McAllester's non-linear square-root bound is avoided due to non-convexity. Instead, we can formally justify our quadratic regularizer by minimizing a linear PAC-Bayesian bound, such as Alquier's bound~\cite{alquier2016properties}. Under Alquier's formulation, assuming a bounded loss function taking values in $[0, 1]$, for any temperature parameter $\lambda' > 0$, with probability at least $1-\delta$ over the random draw of $\mathcal{D}_{\text{cal}}$, the true out-of-distribution multi-task classification risk is bounded by:
\begin{equation}
\label{eq:alquier}
\mathbb{E}_{\tilde{\Theta} \sim Q}[R(G_{\tilde{\Theta}})] \le \mathbb{E}_{\tilde{\Theta} \sim Q}[\widehat{R}_{\text{cal}}(G_{\tilde{\Theta}})] + \frac{D_{\text{KL}}(Q \parallel P) + \ln(1/\delta)}{\lambda'} + \frac{\lambda'}{8 N_{\text{total}}}
\end{equation}
Here, the third term represents the variance contribution from the choice of temperature $\lambda'$ and the dataset size. Note that Alquier's bound strictly assumes a $[0, 1]$-bounded loss function. While the multi-task cross-entropy loss $\ell_{\text{ce}}$ is theoretically unbounded in $[0, \infty)$, we assume that the loss is bounded via a standard clipping threshold $L_{\max}$ (i.e., $\ell_{\text{clipped}} = \min(\ell_{\text{ce}}, L_{\max})$) and rescaled by $1 / L_{\max}$ to reside within $[0, 1]$, which is a standard technical assumption in deep learning generalization bounds~\cite{guedj2019primer}. In our physical PyTorch implementation, this $1 / L_{\max}$ scaling factor is mathematically absorbed into the learning rate and the regularization hyperparameter $\lambda_{\text{PAC}}$, allowing us to minimize the unrescaled clipped loss directly without altering the underlying gradient dynamics. This closes the theory-to-practice gap while maintaining optimal numerical conditioning.

Minimizing this linear PAC-Bayesian bound with respect to the mean parameters $\Theta$ is mathematically equivalent to minimizing the following risk and complexity objective (since the third term of Equation~\eqref{eq:alquier} is a constant independent of $\Theta$):
\begin{equation}
\min_{\Theta} \mathcal{L}_{\text{ce}}(\Theta) + \frac{D_{\text{KL}}(Q \parallel P)}{\lambda'}
\end{equation}
where $\mathcal{L}_{\text{ce}}(\Theta) = \mathbb{E}_{\tilde{\Theta} \sim Q}[\widehat{R}_{\text{cal}}(G_{\tilde{\Theta}})]$. Substituting the analytical form of $D_{\text{KL}}(Q \parallel P)$ from Equation~\eqref{eq:kl_derivation} and dropping terms constant with respect to $\Theta$, the objective reduces exactly to:
\begin{equation}
\min_{\Theta} \mathcal{L}_{\text{total}}(\Theta) = \mathcal{L}_{\text{ce}}(\Theta) + \lambda_{\text{PAC}} \mathcal{R}_{\text{PAC}}(\Theta)
\end{equation}
where the quadratic $L_2$ regularizer is:
\begin{equation}
\label{eq:pac_regularizer}
\mathcal{R}_{\text{PAC}}(\Theta) = \|\Theta - \Theta_{\text{uniform}}\|_2^2 = \sum_{k=0}^{K-1} \left( (\theta_{k,0} - \theta_{\text{uniform}})^2 + \sum_{j=1}^d \theta_{k,j}^2 \right)
\end{equation}
and the regularization trade-off coefficient is theoretically defined as:
\begin{equation}
\lambda_{\text{PAC}} = \frac{1}{2 \sigma_0^2 \lambda'}
\end{equation}
The free parameter $\lambda' > 0$ balances the empirical classification performance and parameter-space complexity. Minimizing the bound in Equation~\eqref{eq:alquier} with respect to $\lambda'$ reveals that the optimal temperature scales as $\lambda'^* = \sqrt{8 N_{\text{total}} (D_{\text{KL}}^* + \ln(1/\delta))}$. In practical deep learning optimization, however, $\lambda'$ is treated as a fixed hyperparameter proportional to the sample size, e.g., $\lambda' = \beta \cdot N_{\text{total}}$, yielding the standard theoretical scaling value $\lambda_{\text{PAC}} = \frac{1}{2 \sigma_0^2 \beta N_{\text{total}}}$. This represents an elegant Lagrangian relaxation of McAllester's complexity bound, where the trade-off coefficient acts as an inverse temperature trade-off balancing empirical risk and parameter-space complexity.
\begin{equation}
\mathcal{L}_{\text{ce}}(\Theta) = \frac{1}{K \cdot M} \sum_{k=0}^{K-1} \sum_{i=1}^{M} \ell_{\text{ce}}\left(f\left(x_i^{(k)}; W_{\text{merged}}(\Theta)\right), y_i^{(k)}\right)
\end{equation}

\begin{remark}
We must explicitly clarify that while the PAC-Bayesian bound provides a mathematically rigorous qualitative regularizer (guaranteeing that the complexity penalty is strictly bounded and theoretically correct), the actual numerical value of the bound under extreme few-shot scarcity ($N_{\text{total}} = 40$) is numerically vacuous (exceeding 1.0 for the error rate), which is a well-known standard property of deep learning generalization bounds. Nonetheless, the derived $L_2$ Consensus-Pulling penalty acts as an exceptionally effective stabilizer during optimization, smoothing out the non-convex loss landscape and preventing transductive overfitting on calibration noise.
\end{remark}

\subsection{Stable Consensus Target Initialization}
\label{subsec:initialization_consensus}
At the beginning of optimization, we desire our ensembling coefficients to represent the stable, scale-preserving uniform ensembling consensus baseline where each expert is merged equally:
\begin{equation}
\alpha_k(l) = \frac{1}{K}, \quad \forall k \in \{0, \dots, K-1\}, \quad \forall l \in \{0, \dots, L-1\}
\end{equation}
To achieve this, we set the prior target parameters $\Theta_{\text{uniform}}$. For a flat, depth-independent initial trajectory, all higher-order polynomial coefficients must be zero: $\theta_{\text{uniform}, j} = 0.0$ for $j \ge 1$. To determine the initial bias parameter $\theta_{\text{uniform}, 0} = \theta_{\text{uniform}}$, we invert the logistic sigmoid mapping:
\begin{equation}
\sigma(\theta_{\text{uniform}}) = \frac{1}{K} \implies \frac{1}{1 + e^{-\theta_{\text{uniform}}}} = \frac{1}{K} \implies 1 + e^{-\theta_{\text{uniform}}} = K
\end{equation}
Solving for $\theta_{\text{uniform}}$ yields:
\begin{equation}
\theta_{\text{uniform}} = -\ln(K - 1) = \ln\left(\frac{1}{K-1}\right)
\end{equation}
For our multi-task setup with $K = 4$ tasks, we compute:
\begin{equation}
\theta_{\text{uniform}} = \ln\left(\frac{1}{3}\right) \approx -1.0986
\end{equation}

By initializing the learnable parameters with $\theta_{k,0} = -1.0986$ and $\theta_{k,j \ge 1} = 0.0$, the initial merged model is mathematically identical to the Static Uniform baseline. Center-pulling the parameters toward $\theta_{\text{uniform}}$ via our $L_2$ penalty mathematically guarantees parameter scale conservation, preventing coordinate-space representation explosion during gradient optimization.

\subsection{SWA Equivalence of Uniform Weight Merging}
\label{subsec:swa_equivalence}
To understand why the Static Uniform baseline serves as an exceptionally robust prior target, we establish a formal theoretical link between weight-space model merging and Stochastic Weight Averaging (SWA)~\cite{izmailov2018averaging}.

\begin{theorem}
Let $W_0 \in \mathbb{R}^D$ be a shared pre-trained base network, and let $W_k = W_0 + \delta_k$ be $K$ task-specific experts fine-tuned on task-specific losses $\mathcal{L}_k(W)$. Suppose that during independent fine-tuning, the expert parameters are corrupted by independent task-specific SGD sampling noise $\epsilon_k \sim \mathcal{N}(0, \sigma^2 I)$ centered at their respective local minima $\theta_k^*$, such that $\delta_k = \theta_k^* - W_0 + \epsilon_k$. Then, the uniform merged model $W_{\text{uniform}} = W_0 + \frac{1}{K} \sum_k V_k$ acts as a parametric low-pass filter that reduces the variance of the optimization noise by a factor of $K$:
\begin{equation}
\mathbb{E}[W_{\text{uniform}}] = W_0 + \frac{1}{K} \sum_{k=0}^{K-1} (\theta_k^* - W_0)
\end{equation}
\begin{equation}
\mathrm{Cov}\left(W_{\text{uniform}}\right) = \frac{\sigma^2}{K} I
\end{equation}
\end{theorem}

\begin{proof}
By definition, the task vectors are $V_k = W_k - W_0 = \theta_k^* - W_0 + \epsilon_k$. The uniform merged model is:
\begin{equation}
W_{\text{uniform}} = W_0 + \frac{1}{K} \sum_{k=0}^{K-1} V_k = W_0 + \frac{1}{K} \sum_{k=0}^{K-1} (\theta_k^* - W_0 + \epsilon_k)
\end{equation}
Taking the expectation with respect to the random SGD noise $\epsilon_k$:
\begin{equation}
\mathbb{E}[W_{\text{uniform}}] = W_0 + \frac{1}{K} \sum_{k=0}^{K-1} (\theta_k^* - W_0) + \frac{1}{K} \sum_{k=0}^{K-1} \mathbb{E}[\epsilon_k] = W_0 + \frac{1}{K} \sum_{k=0}^{K-1} (\theta_k^* - W_0)
\end{equation}
Since the task experts are fine-tuned independently on distinct data streams, their SGD noise vectors $\epsilon_k$ are mutually independent. The covariance is:
\begin{align}
\mathrm{Cov}\left(W_{\text{uniform}}\right) &= \mathrm{Cov}\left(W_0 + \frac{1}{K} \sum_{k=0}^{K-1} (\theta_k^* - W_0 + \epsilon_k)\right) = \mathrm{Cov}\left(\frac{1}{K} \sum_{k=0}^{K-1} \epsilon_k\right) \\
&= \frac{1}{K^2} \sum_{k=0}^{K-1} \mathrm{Cov}(\epsilon_k) = \frac{1}{K^2} \sum_{k=0}^{K-1} \sigma^2 I = \frac{\sigma^2}{K} I
\end{align}
This completes the proof.
\end{proof}

While Theorem 3.1 establishes a useful conceptual link between uniform weight merging and the variance reduction properties of Stochastic Weight Averaging (SWA) under independent sampling noise, its mathematical formulation is intentionally straightforward (leveraging basic central limit covariance properties) and serves primarily as a stylized, conceptual caricature. In practice, its applicability is subject to several crucial modeling limitations. Specifically, modeling disparate task experts fine-tuned on entirely different classification tasks (such as MNIST vs. SVHN) as being corrupted by zero-mean independent SGD noise centered around a single local basin of attraction is highly unrealistic. In real-world multi-task environments, independent fine-tuning drags task expert weights into structurally distinct basins of attraction separated by high non-convex barrier boundaries. This phenomenon is extensively characterized in the Linear Mode Connectivity (LMC) literature~\cite{frankle2020linear, entezari2021role}, which establishes that independently trained networks generally reside in disconnected valleys of the loss landscape. 

Under these conditions of severe multi-task coordinate conflict, uniform ensembling no longer acts as a clean noise canceler; instead, representation conflicts accumulate multiplicatively down successive sequential layers, causing the uniform model to suffer from severe functional collapse. Indeed, the fact that the Static Uniform baseline drops dramatically to \textbf{33.57\%} in our physical 14-layer deep residual network empirically confirms the failure of this single-basin assumption and highlights the breakdown of linear connectivity. This severe baseline degradation mathematically justifies the absolute necessity of our proposed PAC-Bayes Merge framework, which dynamically optimizes layer-wise ensembling coordinates under PAC-Bayesian regularization to navigate distinct basins of attraction, smooth out layer-wise conflicts, and restore multi-task performance.

\subsection{Monte Carlo Expected Risk Optimization and Posterior Ensemble Evaluation}
\label{subsec:monte_carlo}
There is a fundamental mathematical disconnect when one derives a PAC-Bayesian bound on a randomized posterior classifier $G_Q$ but evaluates a single deterministic model compiled at the posterior mean $\Theta^*$. Since deep neural networks are highly non-convex, the risk of the deterministic model at the mean, $R(G_{\mathbb{E}[\tilde{\Theta}]})$, is not mathematically bounded by the randomized expected risk $\mathbb{E}_{\tilde{\Theta} \sim Q}[R(G_{\tilde{\Theta}})]$.

To bridge this theoretical-to-empirical gap, we directly implement the randomized classifier during both optimization and evaluation:
\begin{enumerate}
    \item \textbf{Randomized Training:} At each optimization step, rather than computing empirical risk solely at the mean trajectory parameters $\Theta$, we draw $S$ independent Monte Carlo samples $\tilde{\Theta}_s \sim \mathcal{N}(\Theta, \sigma^2 I)$ and optimize the expected cross-entropy loss over these samples:
    \begin{equation}
    \mathcal{L}_{\text{ce}}(\Theta) \approx \frac{1}{S \cdot K \cdot M} \sum_{s=1}^S \sum_{k=0}^{K-1} \sum_{i=1}^{M} \ell_{\text{ce}}\left(f\left(x_i^{(k)}; W_{\text{merged}}(\tilde{\Theta}_s)\right), y_i^{(k)}\right)
    \end{equation}
    This directly aligns our optimization objective with McAllester's expected risk bound, and acts as a powerful variational regularizer that smooths the loss landscape.
    \item \textbf{Posterior Ensemble Evaluation:} At test time, we evaluate the true randomized classifier $G_Q$ by sampling $S_{\text{test}}$ trajectory coordinates from the posterior distribution and averaging their softmax probability predictions:
    \begin{equation}
    P(y \mid x) \approx \frac{1}{S_{\text{test}}} \sum_{s=1}^{S_{\text{test}}} \mathrm{softmax}\left( f\left(x; W_{\text{merged}}(\tilde{\Theta}_s)\right) \right)
    \end{equation}
    This Monte Carlo approximation exactly represents the randomized PAC-Bayesian classifier, resolving the theoretical mismatch and significantly improving test-time robustness under extreme data scarcity.
\end{enumerate}
